\section{Zero Lelong numbers for nef klt adjoints}
\label{sec:adjoint-metrics}

We prove \Cref{thm:adjoint-metric}. The construction uses sections
only after one fixed twist pulled back from an ample bundle on the
normal base. Its central estimate must reduce a
normalized section's mass without changing its value at a point of
positive Lelong number. We first recall the top-degree $L^2$ estimate,
then construct the normalized extremal sections. The affine localization
lemma turns their small sublevel mass into the required competitor.

\subsection{Weights and the complete \texorpdfstring{$L^2$}{L2} estimate}

A local weight $\varphi$ of a Hermitian metric gives the squared norm
$|\xi|^2e^{-\varphi}$ in its local frame.  Semipositive curvature means
that the local weights are plurisubharmonic.  A metric on a rational
line bundle is defined by taking a positive integral tensor power and
dividing its weights by that integer.  All identifications below are
identifications of rational line bundles, rather than of numerical
classes.

If a weight $\eta$ is a weight on $K_V$, the expression $e^\eta$ denotes
a density: in a canonical coordinate frame it is multiplied by the
corresponding coordinate volume.  Likewise, for an $F$-valued top form
$v$, the expression $|v|^2e^{-\psi}$, with $\psi$ a weight on $F$,
denotes a density.  These densities are invariant under changes of
coordinates and frames.  In particular, the integral of the latter
does not depend on an auxiliary K\"ahler metric.

We shall use the following form of the complete $L^2$ theorem.  The
matrix observation in its statement will be responsible for the
uniform constants.

\begin{lemma}\label[lemma]{lem:top-degree-l2}
Let $(M,\omega_M)$ be a complete K\"ahler manifold of dimension $n$,
and let $F$ be a holomorphic line bundle with a smooth Hermitian metric
of local weight $\psi$ and strictly positive curvature $\theta$.
Write $\Lambda_{\omega_M}$ for the adjoint of wedging with
$\omega_M$.  If $\beta$ is an
$F$-valued, $\dbar$-closed $(n,1)$ form which is square integrable and
satisfies
\[
 I(\beta)=\int_M
 \left\langle[\theta,\Lambda_{\omega_M}]^{-1}\beta,
                 \beta\right\rangle_{\psi,\omega_M}\,dV_{\omega_M}<\infty,
\]
then there is an $F$-valued $(n,0)$ form $v$ with
\[
 \dbar v=\beta,\qquad \int_M |v|^2e^{-\psi}\le I(\beta).
\]
In local coordinates, the inverse-curvature density in $I(\beta)$ is
the inverse Hermitian curvature matrix acting on the remaining
$(0,1)$ covector, multiplied by the coordinate top-form density.  It
is independent of the complete background metric.  Consequently:
\begin{enumerate}[label=\textup{(\roman*)}]
\item if $\theta\ge\omega_0>0$, this density is bounded by the
      $(n,1)$ norm density computed with $\omega_0$;
\item if $g$ is a smooth real function and
      $\theta\ge a\,i\partial g\wedge\bar\partial g$ with $a>0$,
      the inverse-curvature squared norm of $\bar\partial g$ is at most
      $a^{-1}$.
\end{enumerate}
\end{lemma}

\begin{proof}
The existence statement is the complete K\"ahler $L^2$ theorem with
positive curvature operator and finite inverse-curvature integral
\cite[Theorem~5.1, p.~33]{Demailly2012}.  It requires pointwise
positivity and the displayed integral, not a uniform lower spectral
bound relative to $\omega_M$.

For the density calculation, write a top-form-valued $(0,1)$ form as
$\sum_j\beta_j\,dz_1\wedge\cdots\wedge dz_n\wedge d\bar z_j$.
The determinant in the squared norm of the top form cancels the volume
determinant.  Diagonalizing the curvature relative to $\omega_M$
shows that applying the inverse commutator leaves
$\sum_{j,k}(\theta^{-1})^{j\bar k}\beta_j\overline{\beta_k}$
against the coordinate volume and line weight.  All formulas use the
same standard normalization of forms and volume.  Matrix order gives
$\theta^{-1}\le\omega_0^{-1}$ in (i).  For (ii), conjugate the
rank-one inequality $a\,\partial g\otimes\overline{\partial g}
\le\theta$ by $\theta^{-1/2}$.
\end{proof}

\subsection{The fixed twist and normalization}

The zero-dimensional case is
immediate, so put $n=\dim V>0$ and fix a K\"ahler form $\omega$ on
$V$.  Choose compatible canonical divisors
and write
\begin{equation}\label{eq:adjoint-boundary}
 N=K_V+B.
\end{equation}
The support of $B$ is simple normal crossing, every coefficient is
strictly less than one, and the negative part $B^-$ is exceptional over
$H$.  The last assertion follows because the nonexceptional
coefficients are those of the effective boundary $\Theta$.

Let $b=[B]$ be the divisor weight, a rational linear combination of
logarithms of squared absolute values of local divisor equations.  In
crossing coordinates it has the form
\[
 b=\sum_{j=1}^{a}\beta_j\log|z_j|^2+O(1),\qquad \beta_j<1.
\]
Thus $e^{-b}$ is locally integrable.  Negative coefficients cause no
failure of this integrability statement.

Fix a very ample Cartier divisor $H_0$ on $H$, and set
$A_H=(n+1)H_0$ and $A=\pi^*A_H$.  Give $A$ a smooth semipositive
metric with weight $a$.  For every positive integer $m$ for which
$mD_H$ is Cartier, the line bundle $mD_H+A_H$ is globally generated.
Indeed, for $1\le j\le n$,
\[
 mD_H+A_H-jH_0-(K_H+\Theta)
 =(m-1)D_H+(n+1-j)H_0
\]
is ample.  Klt Kawamata--Viehweg vanishing, in the form of
\cite[Theorem~3.2]{Fujino2011}, gives the cohomology vanishings for
Castelnuovo--Mumford regularity.  The global-generation conclusion of
\cite[Theorem~1.8.5(i)]{Lazarsfeld2004I} then applies.  Pullback proves
that $mN+A$ is globally generated as well.

Since $N$ is nef, it is pseudo-effective.  Choose a semipositive metric
$\phi$ on $N$ with minimal singularities.  Such a metric is obtained
as the upper envelope, relative to a smooth reference weight, of the
normalized plurisubharmonic weights; see
\cite[Definition~1.4 and Theorem~1.5]{DPS2001}.  Its defining property is
that every other semipositive weight $\psi$ on $N$ satisfies
\begin{equation}\label{eq:minimal-weight-comparison}
 \psi\le\phi+C_\psi
\end{equation}
for a constant $C_\psi$.  We shall show that $\phi$ has zero Lelong
numbers.

Suppose, to the contrary, that its Lelong number at $p\in V$ is
positive.  For each allowed $m$ define
\begin{equation}\label{eq:extremal-gauge}
 Q_m(s)=\int_V |s|^{2/m}e^{-a/m-b},
 \qquad s\in H^0(V,mN+A).
\end{equation}
The integrand is a canonical density by \Cref{eq:adjoint-boundary}.
The integral is finite, continuous in $s$, positive for $s\ne0$, and
homogeneous of degree $2/m$.  These facts also show that $Q_m(s)=1$ is
a compact subset of the finite-dimensional section space: restrict
$Q_m$ first to the unit sphere of any vector-space norm and use its
positive minimum.  Choose $s_m$ maximizing the absolute evaluation at
$p$ on this compact set, using any fixed norm on the one-dimensional
fiber for this $m$.  Global generation implies
\begin{equation}\label{eq:extremal-value}
 Q_m(s_m)=1,\qquad s_m(p)\ne0.
\end{equation}
Put $\tau_m=m^{-1}\log|s_m|^2$, a weight on $N+A/m$.

\begin{claim}\label{clm:extremal-compactness}
After passage to an unbounded subsequence of allowed integers, the
weights $\tau_m-a/m$ are uniformly locally bounded above and converge
locally in $L^1$ and almost everywhere to a semipositive weight
$\tau_\infty$ on $N$. This compactness assertion holds for every
sequence of sections with $Q_m(s_m)=1$; it does not require their
extremality at $p$.
\end{claim}

\begin{proof}
Choose a smooth reference weight $\phi_{\rm ref}$ on $N$ and write
$u_m=\tau_m-a/m-\phi_{\rm ref}$.  These are global quasi-psh functions
whose complex Hessians have one common lower bound $-C\omega$ on the
compact connected manifold $V$.  Normalize them by subtracting
$M_m=\sup_Vu_m$.  The compactness theorem for sup-normalized quasi-psh
functions \cite[Proposition~2.7]{GuedjZeriahi2005} gives, along a
subsequence, local $L^1$ and almost-everywhere
convergence of $u_m-M_m$ to a quasi-psh function $u_\infty$ that is
finite almost everywhere.  In particular the normalized functions do
not converge identically to $-\infty$.

The fixed density $e^{\phi_{\rm ref}-b}$ is integrable.  Since
$u_m-M_m\le0$, dominated convergence gives
\[
 J_m:=\int_V e^{u_m-M_m}e^{\phi_{\rm ref}-b}
 \longrightarrow
 J_\infty:=\int_V e^{u_\infty}e^{\phi_{\rm ref}-b}\in(0,\infty).
\]
The positivity uses finiteness of $u_\infty$ almost everywhere, not a
pointwise lower bound.  The normalization in
\Cref{eq:extremal-value} says $e^{M_m}J_m=1$.  Thus $M_m$ converges to
a finite real number, proving the asserted upper bounds and
convergence without the sup normalization.  Finally,
$i\partial\bar\partial(\tau_m-a/m)\ge
-m^{-1}i\partial\bar\partial a$; the error tends to zero.  The limit
weight is therefore semipositive.
\end{proof}

Define the globally scalar difference, almost everywhere,
\[
 g_m=\phi+a/m-\tau_m.
\]
Its values on the zero divisor of $s_m$ may be assigned arbitrarily
when taking integrals; that divisor has measure zero.
By \Cref{eq:minimal-weight-comparison},
$\tau_\infty\le\phi+C_\infty$.  Choose once and for all
$R>C_\infty+1$.  The preceding convergence and domination imply
\begin{equation}\label{eq:extremal-tail}
 \varepsilon_m:=\int_{\{g_m<-R\}}e^{\tau_m-a/m-b}
 \longrightarrow0.
\end{equation}
Indeed, the indicators tend to zero almost everywhere, because
$\phi-\tau_\infty\ge-C_\infty$ where both weights are finite; the
densities have one integrable majorant on a fixed finite coordinate
cover.

It is now enough to construct sections $w_m\in H^0(V,mN+A)$ such that
\[
 w_m(p)=s_m(p),\qquad Q_m(w_m)^m\le C_0\varepsilon_m,
\]
where $C_0$ is independent of $m$.  For large $m$ this gives
$Q_m(w_m)<1$, so scalar normalization would produce a section of
gauge one with larger absolute evaluation at $p$.  We next construct
these competitors and prove the two displayed properties.

\subsection{A singular weighted equation with a fixed constant}

We will cut off $s_m$ on $\{g_m<-R\}$ and correct the resulting
$\dbar$ error.  The weight for this equation must serve two purposes.
Its singularity near $p$ must force the correction to vanish at $p$
after the holomorphic competitor has been extended.  Its growth where $g_m$
is positive must also allow a uniform conversion from the quadratic
estimate to $Q_m$.
The first requirement determines its slope $c$ near $-\infty$; a fixed
upper slope bound $d$ will give the second.

Positive Lelong number gives local coordinates at $p$ and a sufficiently
small fixed $\alpha>0$ with
$\phi(z)\le\alpha\log|z|^2+O(1)$.  On a cone of positive angular
measure on which all crossing coordinates have modulus comparable to
$|z|$, the weight $e^{-c\phi-b}$ is bounded below by a positive
constant times
$|z|^{-2(c\alpha+\sum\beta_j)}$.  Choose a fixed $c>0$ large enough
that $c\alpha+\sum\beta_j\ge n$.  Then
\begin{equation}\label{eq:nonintegrable-weight}
 e^{-c\phi-b}\quad\hbox{is not locally integrable at }p.
\end{equation}
This choice is valid even if some $\beta_j$ are negative.

Choose a smooth function $0\le\chi\le1$ equal to one on
$(-\infty,-R-1]$ and to zero on $[-R,\infty)$.  Choose a smooth
convex increasing function $f$ which is affine of slope $c$ near
$-\infty$, has $f''>0$ on the closed transition interval
$[-R-1,-R]$, and satisfies $c\le f'\le d$ for one fixed $d$.  Such a
function is obtained by prescribing a nonnegative compactly supported
second derivative positive on that interval.  In particular there is
a fixed constant $C_f$ with
\begin{equation}\label{eq:convex-weight-bound}
 f(t)\le d\max(0,t)+C_f\qquad(t\in\R).
\end{equation}
All of $R,c,\chi,f,d,C_f$ remain fixed as $m$ varies.

For a fixed $m$, choose a dense smooth affine open $U=U_m\subset V$
avoiding the supports of $B$ and $\operatorname{div}(s_m)$. This open
set need not contain $p$. We will recover the value at $p$ only after
extending the resulting holomorphic section to all of $V$. On $U$
the function $g=g_m$ is psh, since the weight $\tau_m$ is locally
pluriharmonic there and the metrics $\phi,a$ are semipositive.  Put
\begin{equation}\label{eq:local-adjoint-weight}
 F_m=mN+A-K_V,\qquad
 S=(m-1)\tau_m+a/m+b.
\end{equation}
This is a smooth semipositive weight on $F_m|_U$.  Divisor weights and
$\tau_m$ are pluriharmonic on $U$, so its curvature is
$m^{-1}i\partial\bar\partial a$.  We regard $s_m$ as an
$F_m$-valued top form.

\begin{lemma}[Localization of top forms]\label[lemma]{mm:lem:localization}
Let $U$ be a smooth connected affine complex variety of dimension $n\ge1$.
Let $F$ be a holomorphic line bundle on $U$ with a smooth semipositive
metric of weight $S$, let $g$ be a global plurisubharmonic function
not identically $-\infty$, and let $s$ be a holomorphic $F$-valued
$(n,0)$-form. Fix $R\in\R$, a smooth cutoff $0\le\chi\le1$ equal
to one for $t\le-R-1$ and zero for $t\ge-R$, and a smooth convex
function $f$ that is affine of positive slope $c$ near $-\infty$,
with $c\le f'\le d$ for some $d$ and $f''>0$ on $[-R-1,-R]$.
If
\[
 M=\int_{\{g<-R\}}|s|^2e^{-S}<\infty,
\]
there is a measurable $F$-valued top form $u$ such that
\begin{equation}\label{mm:eq:localization-estimate}
 \dbar u=\dbar(\chi(g)s),\qquad
 \int_U|u|^2e^{-S-f(g)}\le C_{\chi,f}M,
\end{equation}
where the equation is distributional and
\begin{equation}\label{mm:eq:localization-constant}
 C_{\chi,f}=\max_{[-R-1,-R]}
 \frac{|\chi'|^2e^{-f}}{f''}.
\end{equation}
The constant is independent of $U,F,S,g,s$ and of the auxiliary
complete K\"ahler metrics. At $g=-\infty$ we set $\chi(g)=1$ and
$f(g)=-\infty$; that locus has measure zero.
\end{lemma}

\begin{proof}
Embed $U$ as a closed complex submanifold of some $\C^a$ and let
$\rho=|z|^2|_U$.  The tubular-neighborhood theorem supplies a
holomorphic retraction of a neighborhood onto $U$
\cite{DocquierGrauert1960}; see also the formulation in
\cite[Theorem~3.1]{Forstneric2010}.  Compose $g$ with the retraction.
Radial convolution in that neighborhood, at radii decreasing to zero,
gives smooth psh functions $g_j\downarrow g$ on common smaller
domains.  More precisely, choose increasing regular exhaustion values
$a_j\to\infty$ and $\Omega_j=\{\rho<a_j\}$.  Use one fixed nonnegative smooth radial averaging kernel of total mass
one. The convolution radius at stage $j$ is at most $1/j$, is small
enough to work near $\overline{\Omega_j}$, and is smaller than all
previous radii.  Radial
averages of a psh function are monotone in the radius, so these choices
give monotonicity wherever two of the functions are defined.  They
also give $g_j\ge g$.

Each $\Omega_j$ has a complete K\"ahler metric dominating
$i\partial\bar\partial\rho$, for example the metric with potential
$\rho-\log(a_j-\rho)$.  For a decreasing positive sequence
$\delta_j\to0$, use the line weight
\[
 P_j=S+f(g_j)+\delta_j\rho.
\]
It is smooth near $\overline{\Omega_j}$, its curvature is strictly
positive, and
\[
 i\partial\bar\partial P_j
 \ge f''(g_j)i\partial g_j\wedge\bar\partial g_j.
\]
The smooth closed datum
$\beta_j=\chi'(g_j)\bar\partial g_j\wedge s$ is square
integrable for the complete metric.  To see this even near the boundary
of $\Omega_j$, use the top-degree density calculation: enlarging the
background metric decreases the $(n,1)$ density, and all coefficients
and line weights are smooth on a neighborhood of the relatively
compact closure.

By \Cref{lem:top-degree-l2}, the inverse-curvature integral is at most
\[
 C\int_{\Omega_j\cap\{g_j<-R\}} |s|^2e^{-S}
 \le C\int_{\{g<-R\}}|s|^2e^{-S},\qquad
 C=\sup_{[-R-1,-R]}\frac{|\chi'|^2e^{-f}}{f''}.
\]
The nonnegative term $\delta_j\rho$ only decreases the integrand.
The last integral is $M$ by definition.
Thus we obtain solutions $u_j$ on $\Omega_j$ with
$\int_{\Omega_j}|u_j|^2e^{-P_j}\le CM$.  The background
complete metrics have introduced no comparison factor.

For completeness, the singular passage uses fixed Hilbert spaces.
If $k\le l$ are fixed and $j\ge l$, then on $\Omega_k$ we have
$P_j\le P_l$, and hence
\[
 \int_{\Omega_k}|u_j|^2e^{-P_l}\le CM.
\]
Take a diagonal weakly convergent subsequence in these countably many
fixed earlier smooth weighted spaces.  Their limits agree as
distributions on overlaps, and define $u$ on $U$.  Weak lower
semicontinuity, then monotone convergence as $l\to\infty$, gives
\[
 \int_{\Omega_k}|u|^2e^{-S-f(g)}\le CM.
\]
Exhausting $U$ proves the required global inequality.  Finally
$\chi(g_j)s\to\chi(g)s$ locally in $L^2$ by bounded convergence,
where $\chi(-\infty)=1$.  Thus the equations pass to distributions
and give \Cref{mm:eq:localization-estimate}.
\end{proof}

We apply the lemma to the adjoint data already fixed on $U$.
The exact density identity is
\begin{equation}\label{eq:normalized-top-density}
 |s_m|^2e^{-S}=e^{\tau_m-a/m-b},
\end{equation}
so its sublevel mass is $M=\varepsilon_m$. Applying
\Cref{mm:lem:localization} with $F=F_m$, $g=g_m$, and $s=s_m$ gives an
$F_m$-valued top form $u$ on $U$ with
\begin{equation}\label{eq:singular-cutoff-estimate}
 \dbar u=\dbar(\chi(g)s_m),\qquad
 \int_U |u|^2e^{-S-f(g)}\le C\varepsilon_m,
\end{equation}
where $C=C_{\chi,f}$ is independent of $m$ and of $U$.

\subsection{Extension, evaluation, and the extremal contradiction}

For $m>d+1$ set
\[
 G=S+d\max(0,g),\qquad w_m=\chi(g)s_m-u\quad\hbox{on }U.
\]
\Cref{eq:singular-cutoff-estimate} shows that $w_m$ is holomorphic.  Since $G\ge S$, the cutoff term has $G$-norm squared at
most $\varepsilon_m$.  Since $S+f(g)\le G+C_f$, the $G$-norm squared
of $u$ is at most $e^{C_f}C\varepsilon_m$.  Therefore
\begin{equation}\label{eq:competitor-norm}
 \int_U |w_m|^2e^{-G}\le C'\varepsilon_m,
 \qquad C'=2+2e^{C_f}C.
\end{equation}
This constant is independent of $m$.

To extend the section, rewrite its weight as
\begin{equation}\label{eq:extension-weight}
 G=b+a/m+(m-1-d)\tau_m
           +d\max\{\tau_m,\phi+a/m\}.
\end{equation}
It is locally bounded above away from $\Supp B^-$: the coefficient
$m-1-d$ is positive, the psh weights in the last two terms are locally
bounded above, and $b$ has that property away from its negative
components.  Consequently \Cref{eq:competitor-norm} gives ordinary
local $L^2$ bounds there.  Holomorphic $L^2$ removal across analytic
sets extends $w_m$ to $V\setminus\Supp B^-$.

There is an open $H^\circ\subset H$, with complement of codimension
at least two, over which $\pi$ is an isomorphism.  It misses the image
of $B^-$.  The section just obtained descends on $H^\circ$ to a
section of the actual Cartier bundle
$\cO_H(mD_H+A_H)$.  Normality extends it across $H\setminus
H^\circ$: after trivializing that line bundle, this is the Hartogs
extension of regular holomorphic functions on a normal space.
Pullback now gives a global section, again denoted $w_m$, of $mN+A$.
By projectivity and GAGA it is an algebraic global section as well.
It agrees with the constructed section on $U$ by the identity theorem.
In particular \Cref{eq:competitor-norm} continues to concern the same
section.  No estimate across the negative exceptional divisor was
needed for this extension.

Near $p$, the weight $\tau_m$ is smooth by
\Cref{eq:extremal-value}, whereas $g\to-\infty$.  Thus $\chi(g)=1$
and $f(g)=cg+\text{constant}$ near $p$ on $U$, so
\[
 u=s_m-w_m,\qquad S+f(g)=b+c\phi+O_m(1).
\]
If $s_m(p)-w_m(p)$ were nonzero, its holomorphic coefficient would
have a positive lower bound on a smaller neighborhood.  The finite
weighted integral in \Cref{eq:singular-cutoff-estimate} would then
contradict \Cref{eq:nonintegrable-weight}.  Removing the analytic
null set $V\setminus U$ does not change this divergence.  This also
applies when $p$ lies on $B^-$, since both sections are now holomorphic
at $p$.  We conclude
\begin{equation}\label{eq:competitor-value}
 w_m(p)=s_m(p)\ne0.
\end{equation}
The bounded factor denoted $O_m(1)$ is used only to test integrability
for a fixed $m$; it never enters \Cref{eq:competitor-norm}.

It remains to convert that quadratic estimate to the original gauge.
H\"older's inequality gives
\begin{equation}\label{eq:gauge-holder-first}
 Q_m(w_m)^m\le C'\varepsilon_m
 \left(\int_U e^{(G-a-mb)/(m-1)}\right)^{m-1}.
\end{equation}
There is no unbounded factor hidden in the second integral.  Put
$\eta=d/(m-1)\in(0,1)$.  Direct expansion yields
\begin{equation}\label{eq:canonical-density-mixture}
 \frac{G-a-mb}{m-1}
 =(1-\eta)(\tau_m-a/m-b)
  +\eta\bigl(\max\{\tau_m-a/m,\phi\}-b\bigr).
\end{equation}
Both exponentials on the right are canonical densities.  The first
has integral one; the second has integral at most $1+I$, where
\[
 I=\int_V e^{\phi-b}<\infty.
\]
Finiteness follows from local upper bounds for $\phi$ and local
integrability of $e^{-b}$ on a fixed finite coordinate cover.
Applying H\"older to \Cref{eq:canonical-density-mixture} gives
\[
 \left(\int_U e^{(G-a-mb)/(m-1)}\right)^{m-1}
 \le (1+I)^{\eta(m-1)}=(1+I)^d.
\]
All integrals are unchanged by the omitted analytic null set.
Together with \Cref{eq:extremal-tail,eq:gauge-holder-first}, this proves
\[
 Q_m(w_m)^m\le C'\varepsilon_m(1+I)^d\longrightarrow0.
\]
For large $m$, therefore, $0<Q_m(w_m)<1$.  Multiply $w_m$ by
$Q_m(w_m)^{-m/2}>1$.  The new section has gauge one and, by
\Cref{eq:competitor-value}, strictly larger absolute evaluation than
$s_m$.  This contradicts the choice of $s_m$, regardless of the rate
at which $s_m(p)$ or $\varepsilon_m$ tends to zero.  Hence $\phi$ has
no positive Lelong number at any point.  This proves
\Cref{thm:adjoint-metric}.


\begin{corollary}\label[corollary]{cor:minimal-multiplier}
Under the hypotheses of \Cref{thm:adjoint-metric}, let $\phi$ be a
minimal weight on the actual rational line bundle $N$. Then
\[
 \mathcal I(t\phi)=\cO_V\qquad\text{for every real }t>0,
\]
where the ideal is defined by local integrability of $|f|^2e^{-t\phi}$.
The conclusion is independent of the choice of minimal metric and of
the Cartier multiple used to normalize its weights.
\end{corollary}
\begin{proof}
Skoda's integrability theorem gives local integrability of
$e^{-t\phi}$ for every fixed $t>0$ when the Lelong numbers vanish
\cite[Proposition~7.1]{Skoda1972}; for this formulation for arbitrary
plurisubharmonic weights, see \cite[Property~1.4(8)]{DemaillyKollar2001}
and \cite[Lemma~5.6(a)]{Demailly2012}. Every holomorphic germ is bounded
on a smaller neighborhood, so the multiplier ideal is the structure
sheaf. Two metrics with minimal singularities have weights differing
by a bounded function, directly from their defining comparison.
Dividing the weight of a positive Cartier multiple by that multiple
preserves the stated normalization. These observations prove the
independence assertions.
\end{proof}
